% ══════════════════════════════════════════════════════════════════════════════
% Préambule du rapport de stage académique (École Centrale Casablanca)
%
% Deux exigences se superposent ici, et il a fallu les tenir ensemble :
%
%   1. le guide de rédaction de l'école — police 12, interligne 1,5, texte
%      justifié, pagination ;
%   2. la charte des livrables du projet — c'est le gabarit des rapports
%      hebdomadaires remis au cabinet (rapports/_preambule.tex) qui est repris
%      ici : même page de garde sur fond navy, mêmes bandeaux de partie à
%      pastille numérotée, même table des matières, mêmes en-têtes.
%
% Ce qui a été adapté : la typographie du corps (12 pt et interligne 1,5 au lieu
% de 11 pt serré), les marges, la page de garde qui doit porter les mentions
% imposées (numéro d'étudiant, tuteur, dates), et un bandeau sans numéro pour
% les sections non numérotées (remerciements, introduction, conclusion…).
% ══════════════════════════════════════════════════════════════════════════════

% ─── Encodage & langue ────────────────────────────────────────────────────────
\usepackage[T1]{fontenc}
\usepackage[utf8]{inputenc}
\usepackage[french]{babel}
% Pas de lmodern : les rapports hebdomadaires composent en Computer Modern (les
% fontes T1 « cm-super »), et le rapport doit être dans la même police qu'eux.

% ─── Mise en page imposée par le guide ────────────────────────────────────────
\usepackage[left=2.5cm, right=2.5cm, top=2.4cm, bottom=2.5cm]{geometry}
\usepackage{setspace}
\onehalfspacing                     % interligne 1,5
\usepackage{ragged2e}
\justifying                         % texte justifié
\setlength{\parindent}{1.5em}    % même retrait d'alinéa que les rapports hebdo
\setlength{\parskip}{0.22em plus 0.08em minus 0.04em}
\usepackage{indentfirst}
\setlength{\headheight}{15pt}
\setlength{\headsep}{18pt}

% ─── Mathématiques, graphiques & tableaux ─────────────────────────────────────
\usepackage{amsmath,amssymb}
\usepackage{graphicx}
\graphicspath{{img/}{img/photos/}{../}}
\usepackage{float}
\usepackage{booktabs}
\usepackage{array}
\usepackage{tabularx}
\usepackage{multirow}
\usepackage{xcolor}
\usepackage{colortbl}
\usepackage[font=small,labelfont={bf,color=eccblue}]{caption}
\usepackage{tikz}
\usetikzlibrary{arrows.meta,decorations.pathmorphing,patterns,calc,fadings,
	positioning,shapes.geometric}
\usepackage[most]{tcolorbox}
\usepackage{enumitem}
\setlist[itemize,1]{label=\textbullet, leftmargin=1.6em, itemsep=0.22em, topsep=0.32em}
\setlist[enumerate,1]{leftmargin=1.8em, itemsep=0.22em, topsep=0.32em}

% ─── Couleurs (charte du projet, identique aux rapports hebdomadaires) ────────
\definecolor{eccblue}{RGB}{0,82,145}
\definecolor{linkblue}{RGB}{0,0,250}
\definecolor{navy}{RGB}{12,25,48}
\definecolor{accent}{RGB}{0,188,170}
\definecolor{smoke}{RGB}{160,170,185}
\definecolor{darkgray}{RGB}{80,80,80}
\definecolor{choubelor}{RGB}{249,204,40}
\definecolor{arelire}{RGB}{176,0,32}

\newcommand{\coverSerif}{\fontfamily{ptm}\selectfont}
\newcommand{\coverSans}{\fontfamily{phv}\selectfont}

% ─── Liens ────────────────────────────────────────────────────────────────────
\usepackage{hyperref}
\usepackage{bookmark}
\hypersetup{colorlinks=true, linkcolor=linkblue, urlcolor=linkblue,
	citecolor=linkblue, filecolor=linkblue}

% ═══ Bandeaux de partie ═══════════════════════════════════════════════════════
% Repris des rapports hebdomadaires : pastille dégradée portant le numéro,
% barre bleue, titre centré en italique gras, filet de séparation. Les
% coordonnées sont exprimées pour une justification de 16 cm (marges 2,5 cm) ;
% le bandeau déborde de 1,1 cm dans la marge de gauche, comme dans le gabarit
% d'origine, pour que la pastille « sorte » du bloc de texte.
\newcounter{manualbookmark}
\makeatletter

\newcommand{\bandeautitre}[1]{%
	\node[text width=13.6cm,align=center,
	font=\bfseries\itshape\fontsize{17.5}{21}\selectfont,color=black]
	at (9.35,-0.92) {#1};
	\draw[eccblue!55,line width=0.50pt] (1.55,-1.80) -- (17.10,-1.80);}

\newcommand{\mainsection}[1]{%
	\clearpage
	\thispagestyle{empty}%
	\refstepcounter{section}%
	\setcounter{subsection}{0}%
	\stepcounter{manualbookmark}%
	\phantomsection%
	\pdfbookmark[1]{\thesection\space #1}{sec-\themanualbookmark}%
	\addcontentsline{toc}{section}{\protect\numberline{\thesection}#1}%
	\markboth{\thesection. #1}{\thesection. #1}%
	\vspace*{0.05cm}%
	\noindent\hspace*{-1.55cm}%
	\begin{tikzpicture}[x=1cm,y=1cm]
		\shade[inner color=eccblue!18,outer color=eccblue!82]
		(0.85,0.05) circle (1.30);
		\node[anchor=west,font=\fontsize{8.5}{10}\selectfont,color=white]
		at (0.16,0.90) {Partie};
		\node[font=\bfseries\fontsize{32}{32}\selectfont,color=white]
		at (0.87,0.02) {\thesection};
		\fill[eccblue!80] (1.62,-0.16) rectangle (17.10,0.06);
		\bandeautitre{#1}
	\end{tikzpicture}\par
	\vspace{0.70cm}%
}

% Même bandeau, sans numéro : remerciements, glossaire, introduction,
% conclusion, sources, annexes. La pastille passe en teinte d'accent et reste
% muette — écrire un numéro là où il n'y en a pas serait un contresens.
\newcommand{\mainsectionstar}[1]{%
	\clearpage
	\thispagestyle{empty}%
	\stepcounter{manualbookmark}%
	\phantomsection%
	\pdfbookmark[1]{#1}{sec-\themanualbookmark}%
	\addcontentsline{toc}{section}{#1}%
	\markboth{#1}{#1}%
	\vspace*{0.05cm}%
	\noindent\hspace*{-1.55cm}%
	\begin{tikzpicture}[x=1cm,y=1cm]
		\shade[inner color=accent!14,outer color=accent!70]
		(0.85,0.05) circle (1.30);
		\fill[white,opacity=0.85] (0.85,0.05) circle (0.20);
		\fill[accent!80] (1.62,-0.16) rectangle (17.10,0.06);
		\bandeautitre{#1}
	\end{tikzpicture}\par
	\vspace{0.70cm}%
}
\makeatother

% ─── Titres de niveau inférieur ───────────────────────────────────────────────
\usepackage{titlesec}
\titleformat{\subsection}
{\color{eccblue}\large\bfseries}{\thesubsection.}{0.55em}{}
\titleformat{\subsubsection}
{\color{navy}\normalsize\bfseries\itshape}{\thesubsubsection.}{0.5em}{}
\titlespacing*{\subsection}{0pt}{1.25em}{0.45em}
\titlespacing*{\subsubsection}{0pt}{0.95em}{0.35em}

% ─── En-têtes et pieds de page (gabarit hebdomadaire) ─────────────────────────
\usepackage{fancyhdr}
\pagestyle{fancy}
\fancyhf{}
\fancyhead[L]{\scriptsize\leftmark}
\fancyhead[R]{\scriptsize\thepage}
\fancyfoot[L]{\fontsize{8}{9}\selectfont\color{darkgray}\textit{\nouncleve}}
\fancyfoot[C]{\fontsize{9}{10}\selectfont\bfseries\color{eccblue}\thepage}
\fancyfoot[R]{\fontsize{8}{9}\selectfont\color{darkgray}\textit{Stage opérateur --- \entreprise}}
\renewcommand{\headrulewidth}{1.0pt}
\renewcommand{\footrulewidth}{0.8pt}

% ─── Table des matières ───────────────────────────────────────────────────────
\usepackage{tocloft}
\renewcommand{\contentsname}{Table des matières}
\renewcommand{\cfttoctitlefont}{\hfill\LARGE\bfseries\itshape\color{black}}
\renewcommand{\cftaftertoctitle}{\hfill\par\vspace{0.85cm}}
\renewcommand{\cftsecfont}{\small\color{eccblue}\bfseries}
\renewcommand{\cftsecpagefont}{\small\bfseries\color{black}}
\renewcommand{\cftsubsecfont}{\small\color{navy}}
\renewcommand{\cftdotsep}{2.5}

\makeatletter
\newcommand{\printstyledtoc}{%
	\clearpage
	\pagestyle{empty}%
	\thispagestyle{empty}%
	\phantomsection%
	\pdfbookmark[1]{Table des matières}{toc}%
	\vspace*{0.58cm}%
	\noindent
	\begin{tikzpicture}[baseline]
		\fill[gray!28] (0,0.09) rectangle (0.62\linewidth,0.50);
		\node[anchor=east,font=\LARGE\bfseries\itshape,color=black]
		at (\linewidth,0.30) {Table des matières};
		\draw[black,line width=0.55pt] (0,0.00) -- (\linewidth,0.00);
	\end{tikzpicture}\par
	\vspace{0.82cm}%
	{\hypersetup{linkcolor=linkblue,linktoc=section}\@starttoc{toc}}%
	\clearpage
	\pagestyle{fancy}}%
\makeatother

% ─── Encadrés ─────────────────────────────────────────────────────────────────
\newtcolorbox{encadre}[1][]{colback=eccblue!5, colframe=eccblue!55,
	boxrule=0.7pt, arc=2pt, left=9pt, right=9pt, top=7pt, bottom=7pt, #1}
\newtcolorbox{encadrecle}[1][]{colback=accent!7, colframe=accent!70,
	boxrule=0.7pt, arc=2pt, left=9pt, right=9pt, top=7pt, bottom=7pt, #1}
\newtcolorbox{citationbox}[1][]{colback=navy, colframe=navy, coltext=white,
	boxrule=0pt, arc=2pt, left=11pt, right=11pt, top=9pt, bottom=9pt, #1}

% ─── Marqueur des informations restant à compléter ────────────────────────────
% Rend visible, à la relecture, tout ce que seul l'auteur ou le tuteur peut
% renseigner. À supprimer avant l'envoi du rapport.
\newcommand{\acompleter}[1]{{\color{arelire}\bfseries[à compléter : #1]}}

% ─── Insertion d'images tolérante ─────────────────────────────────────────────
\newcommand{\safelogo}[2]{%
	\IfFileExists{#1}{\includegraphics[height=#2,keepaspectratio]{#1}}{\phantom{\rule{0pt}{#2}}}}

% ─── Figure standard, avec crédit optionnel ───────────────────────────────────
% #1 largeur · #2 fichier · #3 légende · #4 crédit (peut être vide)
% La hauteur est bornée : une capture d'écran de page longue dépasserait sinon
% la hauteur du bloc de texte, et LaTeX la repousserait en laissant une page à
% moitié vide.
\newcommand{\figurecredit}[4]{%
	\begin{figure}[H]
		\centering
		\includegraphics[width=#1\linewidth,height=0.74\textheight,keepaspectratio]{#2}
		\caption{#3\ifx\relax#4\relax\else\\{\footnotesize\color{darkgray}\textit{#4}}\fi}
	\end{figure}}

% ═══ Page de garde ════════════════════════════════════════════════════════════
% Reprise de \couverture des rapports hebdomadaires : fond navy pleine page,
% dégradé haut, logos aux deux angles, pastille de titre, titre serif blanc,
% sous-titre en teinte d'accent, blocs d'information étiquetés, ondes en pied
% de page et filet de bas de page. S'y ajoutent les mentions imposées par
% l'école : numéro d'étudiant, tuteur, dates et année académique.
% #1 badge · #2 titre · #3 sous-titre
\newcommand{\couverturestage}[3]{%
	\hypersetup{pageanchor=false}%
	\begin{titlepage}
		\thispagestyle{empty}
		\begin{tikzpicture}[remember picture,overlay]
			\fill[navy] (current page.south west) rectangle (current page.north east);
			\fill[navy,opacity=0.35]
			(current page.north west)
			rectangle ([yshift=14.8cm]current page.south west -| current page.east);
			\fill[navy,path fading=north]
			([yshift=14.8cm]current page.south west)
			rectangle ([yshift=21cm]current page.south west -| current page.east);

			% Logos : école à gauche, organisme d'accueil à droite
			\node[anchor=north west, inner sep=0pt, fill=white, rounded corners=3pt,
			inner xsep=7pt, inner ysep=7pt]
			at ([xshift=2.2cm,yshift=-1.9cm]current page.north west)
			{\safelogo{../logo.jpg}{2.2cm}};
			\node[anchor=north east, inner sep=0pt, fill=white, rounded corners=3pt,
			inner xsep=7pt, inner ysep=7pt]
			at ([xshift=-2.2cm,yshift=-1.9cm]current page.north east)
			{\safelogo{../app/static/img/logo-choubel.jpg}{2.2cm}};

			% Année académique
			\node[anchor=north east, font=\coverSans\fontsize{8.5}{11}\selectfont, text=smoke]
			at ([xshift=-2.2cm,yshift=-5.1cm]current page.north east)
			{Année académique \anneeacad};

			% Pastille de titre
			\node[anchor=south west,
			font=\coverSerif\fontsize{10.5}{12}\selectfont\bfseries,
			text=white, fill=eccblue, fill opacity=0.85, text opacity=1,
			inner xsep=6pt, inner ysep=4pt, rounded corners=2pt]
			at ([xshift=2.2cm,yshift=19.9cm]current page.south west)
			{#1};

			% Titre et sous-titre
			\node[anchor=south west, text width=16.6cm,
			font=\coverSerif\bfseries\fontsize{24}{29}\selectfont, text=white, inner sep=0pt]
			at ([xshift=2.2cm,yshift=14.9cm]current page.south west)
			{\raggedright\hyphenpenalty=10000\exhyphenpenalty=10000 #2};
			\node[anchor=south west, text width=16.6cm,
			font=\coverSerif\fontsize{14}{18}\selectfont, text=accent, inner sep=0pt]
			at ([xshift=2.2cm,yshift=13.1cm]current page.south west)
			{\raggedright #3};
			\draw[accent!60, line width=1.1pt]
			([xshift=2.2cm,yshift=12.5cm]current page.south west)
			-- ([xshift=6.2cm,yshift=12.5cm]current page.south west);

			% Colonne de gauche : élève, puis période
			\node[anchor=north west, font=\coverSans\fontsize{7.5}{9.5}\selectfont, text=accent]
			at ([xshift=2.2cm,yshift=11.7cm]current page.south west) {ÉLÈVE INGÉNIEUR};
			\node[anchor=north west, text width=8cm,
			font=\coverSerif\fontsize{12.5}{16}\selectfont, text=white, inner sep=0pt]
			at ([xshift=2.2cm,yshift=11.0cm]current page.south west)
			{\textbf{HOUNKPETOHOU} \prenomeleve \\[2pt]
				{\fontsize{9.5}{12}\selectfont\color{smoke}
					N\textsuperscript{o} d'étudiant : \numetudiant}};

			\node[anchor=north west, font=\coverSans\fontsize{7.5}{9.5}\selectfont, text=accent]
			at ([xshift=2.2cm,yshift=8.5cm]current page.south west) {PÉRIODE DE STAGE};
			\node[anchor=north west, text width=8cm,
			font=\coverSerif\fontsize{11.5}{15}\selectfont, text=white, inner sep=0pt]
			at ([xshift=2.2cm,yshift=7.8cm]current page.south west)
			{Du \datedebut{} au \datefin \\[2pt]
				{\fontsize{9.5}{12}\selectfont\color{smoke} Soit six semaines --- \villestage}};

			% Filet vertical, comme sur la couverture des rapports hebdomadaires
			\draw[smoke!30,line width=0.4pt]
			([xshift=10.4cm,yshift=11.9cm]current page.south west)
			-- ([xshift=10.4cm,yshift=6.9cm]current page.south west);

			% Colonne de droite : organisme d'accueil, puis tuteur
			\node[anchor=north west, font=\coverSans\fontsize{7.5}{9.5}\selectfont, text=accent]
			at ([xshift=11.2cm,yshift=11.7cm]current page.south west) {ORGANISME D'ACCUEIL};
			\node[anchor=north west, text width=6.4cm,
			font=\coverSerif\fontsize{12.5}{16}\selectfont, text=white, inner sep=0pt]
			at ([xshift=11.2cm,yshift=11.0cm]current page.south west)
			{\textbf{\entreprise} \\[2pt]
				{\fontsize{9.5}{12}\selectfont\color{smoke} \villestage}};

			\node[anchor=north west, font=\coverSans\fontsize{7.5}{9.5}\selectfont, text=accent]
			at ([xshift=11.2cm,yshift=8.5cm]current page.south west) {TUTEUR ORGANISME D'ACCUEIL};
			\node[anchor=north west, text width=6.4cm,
			font=\coverSerif\fontsize{11.5}{15}\selectfont, text=white, inner sep=0pt]
			at ([xshift=11.2cm,yshift=7.8cm]current page.south west)
			{\tuteurnom \\[2pt]
				{\fontsize{9.5}{12}\selectfont\color{smoke} \tuteurfonction}};

			% Ondes de pied de page : le motif des couvertures du projet
			\begin{scope}[shift={([xshift=10.5cm,yshift=3.4cm]current page.south west)},
				accent,line width=0.7pt,opacity=0.22]
				\foreach \a/\p in {0.6/0, 0.45/40, 0.75/80, 0.35/120} {
					\draw plot[domain=-4.5:4.5,samples=80,smooth]
					(\x,{\a*sin(200*\x+\p)});
				}
			\end{scope}

			\draw[smoke!25,line width=0.3pt]
			([xshift=2.2cm,yshift=1.8cm]current page.south west)
			-- ([xshift=-2.2cm,yshift=1.8cm]current page.south east);
			\node[anchor=south west, font=\coverSans\fontsize{7.5}{9}\selectfont, text=smoke!60]
			at ([xshift=2.2cm,yshift=1.0cm]current page.south west)
			{École Centrale Casablanca --- Rapport de stage opérateur de première année
				--- Projet mené en binôme avec \binome};
		\end{tikzpicture}
	\end{titlepage}
	\hypersetup{pageanchor=true}%
}
